% Part: incompleteness
% Chapter: incompleteness-provability
% Section: lob-thm

\documentclass[../../../include/open-logic-section]{subfiles}

\begin{document}

\olfileid[pt]{inc}{inp}{tar}

\olsection{A indefinibilidade da verdade}

A noção de \emph{definibilidade} depende de dispormos de uma semântica
formal para a linguagem da aritmética. Descrevemos um conjunto de fórmulas e
sentenças na linguagem da aritmética. A ``interpretação pretendida''
consiste em ler essas sentenças como afirmações sobre os números naturais, e
tal afirmação pode ser verdadeira ou falsa. Seja $\Struct{N}$ a
!!{structure} com domínio $\Nat$ e a interpretação padrão dos símbolos da
linguagem da aritmética. Então $\Sat{N}{!A}$ significa ``$!A$ é verdadeira
na interpretação padrão''.

\begin{defn}
Uma relação $R(x_1,\dots,x_k)$ de números naturais é \emph{definível} em
$\Struct{N}$ se, e somente se, existe uma fórmula $!A(x_1,\dots,x_k)$ na
linguagem da aritmética tal que, para quaisquer $n_1,\dots,n_k$,
$R(n_1,\dots,n_k)$ se, e somente se,
$\Sat{N}{!A(\num n_1,\dots,\num n_k)}$.
\end{defn}

Em outras palavras, uma relação é definível em $\Struct{N}$ se, e somente
se, é representável na teoria $\Th{TA}$, onde $\Th{TA} =
\Setabs{!A}{\Sat{N}{!A}}$ é o conjunto das sentenças verdadeiras da
aritmética. (Se isso não lhe parece imediatamente claro, retome as
definições e convença-se de que esse é o caso.)

\begin{lem}
Toda relação computável é definível em~$\Struct{N}$.
\end{lem}

\begin{proof}
É fácil verificar que a fórmula que representa uma relação em $\Th{Q}$
define a mesma relação em $\Struct{N}$.
\end{proof}

Podemos agora perguntar se a recíproca também é verdadeira. Isto é, toda
relação definível em~$\Struct{N}$ é computável? A resposta é não. Por
exemplo:

\begin{lem}
A relação de parada é definível em $\Struct{N}$.
\end{lem}

\begin{proof}
Lembre-se de que o teorema da forma normal de Kleene afirma que toda função
computável parcial~$f$ tem um índice~$e$ tal que $f(x) =
U(\umin{s}{T(e,x,s)})$ para todo $x \in \Nat$, onde $U$ e $T$ são
recursivas primitivas e, portanto, totais. Assim, $f(x)$ está definida (O.s.\@, a computação para) sse existe um~$s$ tal que $T(e,x,s)$ vale.

Seja agora $H$ a relação de parada, O.s.\@,
\[
H = \Setabs{\tuple{e,x}}{\lexists[s][T(e, x, s)]}.
\]
Seja $!D_T$ uma fórmula que define $T$ em $\Struct{N}$. Então
\[
H = \Setabs{\tuple{e,x}}{\Sat{N}{\lexists[s][!D_T(\num e, \num x, s)]}},
\]
de modo que $\lexists[s][!D_T(z, x, s)]$ define~$H$ em $\Struct{N}$.
\end{proof}

\begin{prob}
Mostre que $Q(n) \defiff n \in \Setabs{\Gn{!A}}{\Th{Q} \Proves !A}$ é
  definível na aritmética.
\end{prob}

E quanto à própria $\Th{TA}$? Ela é definível na aritmética? Isto é: o
conjunto $\Setabs{\Gn{!A}}{\Sat{N}{!A}}$ é definível na aritmética? O
teorema de Tarski responde negativamente.

\begin{thm}
\ollabel{thm:tarski}
O conjunto das !!{sentence}s verdadeiras da aritmética não é definível na
aritmética.
\end{thm}

\begin{proof} 
Suponha que $!D(x)$ o definisse, O.s.\@, $\Sat{N}{!A}$ sse
$\Sat{N}{!D(\gn{!A})}$. Pelo lema do ponto fixo, existe uma fórmula $!A$
tal que $\Th{Q} \Proves !A \liff \lnot !D(\gn{!A})$ e, portanto,
$\Sat{N}{!A \liff \lnot !D(\gn{!A})}$. Mas então $\Sat{N}{!A}$ se, e
somente se, $\Sat{N}{\lnot !D(\gn{!A})}$, o que contradiz o fato de que
$!D(y)$ deveria definir o conjunto das afirmações verdadeiras da aritmética.
\end{proof}

Tarski aplicou essa análise a uma noção filosófica mais geral de verdade.
Dada qualquer linguagem $L$, Tarski argumentou que uma noção adequada de
verdade para $L$ teria de satisfazer, para cada sentença $X$,
\begin{quote}
`$X$' é verdadeira se, e somente se, $X$.
\end{quote}
O exemplo frequentemente citado de Tarski, para o inglês, é a sentença
\begin{quote}
`A neve é branca' é verdadeira se, e somente se, a neve é branca.
\end{quote}
Contudo, para qualquer linguagem forte o suficiente para representar a
função diagonal e qualquer predicado linguístico $T(x)$, podemos construir
uma sentença $X$ que satisfaça ``$X$ se, e somente se, não
$T(\text{`$X$'})$.'' Como não queremos que um predicado de verdade declare
algumas sentenças simultaneamente verdadeiras e falsas, Tarski concluiu que
não se pode especificar um predicado de verdade para todas as sentenças de
uma linguagem sem, de algum modo, sair dos limites da própria linguagem. Em
outras palavras, o predicado de verdade de uma linguagem não pode ser
definido na própria linguagem.

\end{document}

